\documentclass[12pt,a4paper]{article}
\usepackage[T2A,T1]{fontenc}
\usepackage[utf8]{inputenc}
\usepackage[english,russian]{babel}
\usepackage{XCharter}
\usepackage{amsmath,amssymb,mathtools}
\usepackage{array,longtable,booktabs}
\usepackage[margin=24mm]{geometry}
\usepackage[hidelinks]{hyperref}
\urlstyle{same}
\setlength{\parindent}{0pt}
\setlength{\parskip}{0.45em}
\setlength{\LTleft}{0pt}
\setlength{\LTright}{0pt}
\emergencystretch=2em
\setlength{\tabcolsep}{4pt}
\newcolumntype{P}[1]{>{\raggedright\arraybackslash}p{#1}}
\title{Ревизия корпуса работ Н. П. Брусенцова}
\author{}
\date{}
\begin{document}
\maketitle

\emph{Итог проверки источников, формул и описаний в исследовательском корпусе. Дата обновления: 23.11.2026.}

\section{Краткий итог}

Основное прочтение модели подтверждается первичными материалами: Н. П. Брусенцов описывает отношения между терминами и классами предметов, использует распределение четырёх комбинаций признаков и вводит дополнительное условие сосуществования термина и его противоположности. Это не та же формальная задача, что оценка материальной импликации между двумя целыми высказываниями.

При ревизии найдено одно важное уточнение к прежней формулировке: \textbf{Н. П. Брусенцов сам называет \(\sigma\) третьим значением истинности классового выражения.} В главе 2005 года при переходе к знаковой шкале привходящий статус кодируется числом \(0\), а в векторе \(+-0+\) этот статус стоит в конкретной позиции \(x'y\). Это не таблица истинности произвольного условного высказывания \(P\to Q\). Значение \(0\) функции \(A(x,y)\) в статье 2003 года, напротив, означает антиприсущность. Коды разных работ нельзя отождествлять по одному числу.

Проверка не даёт основания считать, что предложенная система уже доказала общую несостоятельность классической логики или историческую тождественность аристотелевской силлогистике. Для этих выводов нужны дополнительные семантические и исторические аргументы.

\subsection*{Позиция Н. П. Брусенцова: тезис и уточнения}

Н. П. Брусенцов в работе 2010 года различает материальную импликацию и необходимое следование между терминами; прямые формулировки и точные страницы приведены в основном конспекте и строке 1 чек-листа ниже (PDF с. 1–2; печатные с. 88–89). Его объяснение «содержательности» через сущность термина — авторское философское основание, не независимо доказанный факт о причинной или физической связи (Алгебраическая реконструкция, 2008, PDF с. 1). Отдельно, в главе 2005 года, он критикует строгую импликацию Льюиса: исключение \(xy'\) недостаточно без условий существования \(x\) и \(y'\) (PDF с. 1–2). Это самостоятельный аргумент, который нужно отличать от его тезиса о материальной импликации. Проверка устанавливает соответствие пересказа отмеченным источникам и исследует формальные следствия внутри принятых определений; она не устанавливает несостоятельность всей классической логики или историческую тождественность модели Аристотелю.

\textbf{Обозначения этой ревизии.} \(x,y\) — термины/признаки; \(x',y'\) — противоположные термины; \(xy,xy',x'y,x'y'\) — четыре сочетания в таком порядке. \(P,Q\) — целые высказывания; \(\Gamma\) — множество посылок; \(\Gamma\models Q\) означает семантическое следование заключения, а \(\Gamma\vdash Q\) — существование вывода в выбранном исчислении. \(x\Rightarrow y\) у Н. П. Брусенцова относится к терминам, а не обозначает ни одну из этих двух мета-логических записей. \(X,Y\) — классы носителей признаков \(x,y\), а \(t\) — отдельный предмет; \(X(t)\), \(Y(t)\) означают, что он имеет соответствующий признак. В статье 2010 года \(Axy\) — общеутвердительное суждение, тогда как \(A(x,y)\) в статье 2003 года обозначает статусную функцию. \(x\to y\) — импликация в алгебре терминов; \(P\to Q\) относится к целым высказываниям. В формулах преобразования \(x\to y\equiv(x'\lor y)\equiv\neg(xy')\) символы \(\lor\) и \(\neg\) применены к терминам/классовым выражениям. \(V,V'\) — специальные операции модели (в статье 2001 года \(V\) назван интегральной дизъюнкцией, а \(V'\) отмечает антисуществование соответствующей ячейки), не обычные \(\lor,\neg\). \(I,I'\) — обозначения статьи 2010 года с неполной локальной семантикой; \(\equiv\) в этих цепочках выражает авторское равенство/тождество, но точное общее правило для знака в статье отдельно не задано.

\begin{itemize}
\item В статье 2003 года \(A(x,y)\) — функция статуса отношения: \(1\) означает необходимую присущность, \(0\) — антиприсущность, \(\sigma\) — привходящий статус. Условие \(0<\sigma<1\) (сигма больше нуля и меньше единицы) задаёт положение символа между крайними значениями, но конкретное число не указано; источник не уточняет, допускаются ли разные числовые реализации третьего статуса.
\item В главе 2005 года знаки \(+1,-1,0\) означают соответственно необходимое включение, необходимое исключение и привходящий статус. В профиле \(+-0+\) четыре знака относятся к \(xy,xy',x'y,x'y'\); ноль стоит на \(x'y\). В обобщённой форме главы та же ячейка помечена \(\sigma\): это соответствие между двумя представлениями статуса внутри изложения 2005 года, а не равенство \(\sigma=0\) и не подстановка \(\sigma\) в функцию \(A(x,y)\) 2003 года.
\item \(+-++\) — наш пересчёт истинности точечной материальной импликации \(X(t)\to Y(t)\) по четырём сочетаниям принадлежности. Здесь \(+\) означает истинность, а \(-\) — ложность; в авторском \(+-0+\) знаки обозначают статусы сочетаний классов. Это разные типы оценок.
\item Слова «необходимое включение/исключение» обозначают статусы в алгебраической модели Н. П. Брусенцова; сами по себе они не устанавливают метафизическую, физическую или причинную необходимость.
\item \(Axy\) в статье 2010 года — общеутвердительное суждение «Все \(x\) суть \(y\)»; не следует смешивать его по одной букве \(A\) с функцией \(A(x,y)\) 2003 года. \(V,V'\) — специальные операции модели (в работе 2001 года \(V\) назван интегральной дизъюнкцией, \(V'\) отмечает антисуществование соответствующей ячейки), \(I,I'\) — обозначения статьи 2010 года. \(V\) не равно обычному \(\lor\). В этой формуле нижний индекс после \(V\) называет термин, а штрих у термина — его противоположность. В формуле условия универсума \(V_xV_{x'}V_yV_{y'}\equiv1\) знак \(1\) означает единичное значение авторской алгебры, не количество предметов. Полное формальное толкование \(I\), соседнего соединения факторов и \(\equiv\) в статье 2010 года остаётся ограниченным.
\end{itemize}

\section{Как проходила ревизия}

Работа шла параллельно по трём независимым направлениям:

\begin{enumerate}
\item Формально-источниковый рецензент сопоставил карту работ и основной конспект с ключевыми PDF по силлогистике и следованию: 2003, 2005, 2007–2008, 2010–2011 и 2014 годы.
\item Специалист по Lean отдельно сверил электронные страницы 2–5 главы 2005 года по изображениям PDF и проверил, какие части модели можно доказательно формализовать.
\item Библиографический рецензент сверил записи с указателем МГУ и проверил доступность электронных копий. Результат — 207 записей с явной оговоркой, что полнота установлена относительно списка МГУ, а не всего возможного наследия автора.
\end{enumerate}
Итоговые выводы ниже повторно сверены с самими изображениями ключевых страниц. Все 18 локальных PDF были учтены в карте корпуса, но это \textbf{не} означает, что каждая формула каждой страницы всех 18 работ прошла полную визуальную транскрипцию.

\section{Чек-лист соответствия ключевых утверждений источникам}

В таблице приведены ключевые утверждения проекта, для каждого указаны проверяемый первоисточник, точное место и вид основания. Это выборочная проверка опорных тезисов, а не построчная экспертиза каждой фразы и формулы всего корпуса. «Подтверждено» означает, что формулировка проекта соответствует тому, что именно говорится в указанном месте; это \textbf{не} означает, что авторский исторический или философский вывод доказан. Цитаты и буквальный перевод приведены в основном конспекте, чтобы не пересказывать их по памяти. Номера печатной страницы и страницы PDF разделены; если печатное соответствие не установлено, приведена только страница PDF.

\begin{longtable}{P{0.13\textwidth}P{0.39\textwidth}P{0.15\textwidth}P{0.25\textwidth}}
\toprule
\textbf{Тезис} & \textbf{Источник и основание проверки} & \textbf{Тип тезиса} & \textbf{Результат проверки} \\
\midrule
\endfirsthead
\toprule
\textbf{Тезис} & \textbf{Источник и основание проверки} & \textbf{Тип тезиса} & \textbf{Результат проверки} \\
\midrule
\endhead
Материальная импликация не тождественна необходимому следованию; в УА автор утверждает их равнозначность. Основной конспект, саммари, README. & Н. П. Брусенцов, «Воссоздание…» (2010): PDF с. 1–2, печатные с. 88–89. Прямая формулировка и локаторы приведены в основном конспекте; страницы PDF просмотрены. & Прямой тезис автора. & Подтверждено с оговоркой об УА. Без неё пересказ был бы сильнее источника. \\
Содержательность \(x\Rightarrow y\) автор объясняет через сущность \(x\)-вещей; это не независимое доказательство физической или причинной связи. Основной конспект, саммари. & Н. П. Брусенцов, «Алгебраическая реконструкция силлогистики» (2008), PDF с. 1. Короткая прямая цитата дана в основном конспекте; печатное соответствие PDF-страницы не установлено. & Авторское философское объяснение. & Подтверждено как объяснение автора. В проекте не представлено как установленный факт о сущностях. \\
Аргумент о строгой импликации Льюиса отделён от критики материальной импликации. Основной конспект, карта, Lean. & Н. П. Брусенцов, «Реанимация аристотелевой силлогистики» (2005), PDF с. 1–2, печатные с. 140–141. Фраза о недостаточности одного исключения \(xy'\) и формулы \(V_xV'_{xy'}V_{y'}\), \(V_{xy}V'_{xy'}V_{x'y'}\) сверены по изображениям. & Авторский аргумент и формальная запись. & Подтверждено. Не называется определением материальной импликации. \\
Профиль \(+-0+\) имеет порядок \(xy,xy',x'y,x'y'\); ноль в позиции \(x'y\) — привходящий статус. Основной конспект, саммари, руководство формул. & Н. П. Брусенцов, «Реанимация…» (2005), PDF с. 2–3, печатные с. 141–142. Авторское выражение «привходящим статусом» и вектор сверены по изображениям. & Авторское определение статусов. & Подтверждено. Профиль не объявлен таблицей истинности двух целых высказываний. \\
Операции и статусы в проверенных работах & В 2003 году показаны отрицание-дополнение, инверсия терминов/предикатов и алгебраические преобразования силлогизмов (производная PDF, с. 5–6); в 2005 году на знаковой шкале заданы min/max и смена знака (PDF с. 2, печатная с. 141); в статье 2010 года отношение имеет значения «необходимо / возможно / невозможно» (PDF с. 1–2, печатные с. 86–87). Формулы просмотрены на изображениях. & Источник задаёт операции для своих терминов, классовых выражений и статусной шкалы. & В проверенных местах нет полной таблицы связок для произвольных целых высказываний со значением \(\sigma\); перенос требует отдельной формализации. \\
Критика закона исключённого третьего & В 2003 году закон назван хрисипповым постулатом, несправедливо приписываемым Аристотелю (PDF с. 1); в статье 2010 года сказано, что стоики ввели его против трёхзначности отношения следования (PDF с. 1). & Прямая атрибуция позиции автора и отдельный исторический тезис. & Источники подтверждают, что автор так утверждает; они не дают полного пропозиционального исчисления, из которого исторический тезис следовал бы формально. \\
Функция \(A(x,y)\) 2003 года и знаковый ноль профиля 2005 года не используют один и тот же ноль. Основной конспект, саммари, Lean. & «Трехзначная интерпретация…» (2003), производная PDF-копия, PDF с. 1–2: \(A(x,y)=0\) — антиприсущность; \(\sigma\) «не присуще с необходимостью и не антиприсуще», \(0<\sigma<1\). Статусы 2005 года сверены отдельно по PDF с. 2. & Сопоставление двух авторских кодировок. & Подтверждено. Печатная пагинация преобразованной копии не установлена; число \(\sigma\) источник не задаёт. \\
Формула 2010 года транскрибирована как \((x\Rightarrow y)\equiv Axy\equiv I_{xy}I'_{xy'}I_{x'y'}\); правила соседнего соединения факторов полностью не определены в статье. Руководство формул, основной конспект. & «Воссоздание…» (2010), PDF с. 1–2, печатные с. 88–89. Штрихи и последний множитель \(I_{x'y'}\) сверены по изображению; текстовые извлечения использованы только для поиска. & Транскрипция формулы плюс оценка полноты её определения. & Формула подтверждена визуально. Недоопределённость семантики факторов прямо отмечена; тождество \(I\)- и \(V\)-операций не приписывается автору. \\
Оператор \(V\) — специальная интегральная дизъюнкция; \(V'_{xy'}\) в записи отношения связано с исключением/антисуществованием ячейки. Руководство формул, ревизия. & Н. П. Брусенцов, «Трехзначная диалектическая логика» (2001), опубликованный HTML-текст, разделы с формулами \(V_xV_{x'}\equiv1\) и \(V_{xy}V'_{xy'}V_{x'y'}\); автор называет \(V\) «интегральной дизъюнкцией». В HTML нет печатных разрывов страниц; библиографический диапазон публикации — с. 36–44. См. \href{https://www.computer-museum.ru/histussr/trilog11.htm}{полный текст}. & Авторское определение операции и её контекстное чтение. & Подтверждено по тексту публикации. Точную печатную страницу формулы не удалось установить; ссылка ведёт на полный текст. \\
Статья 2006 года называет \(xy\lor\sigma x'y\lor x'y'\) аналогом материальной импликации с устранёнными парадоксами. Основной конспект, карта. & Н. П. Брусенцов, «Содержательное следование, сосуществование противоположностей и формальные импликации» (2006), PDF с. 1, печатная с. 342. Формулировка цитируется в основном конспекте; подробное объяснение ячейки \(x'y\) проверено отдельно по главе 2005 года. & Краткая авторская характеристика. & Подтверждено. Подробный аргумент 2005 года не приписывается статье 2006 года. \\
Общее суждение не даёт частного вывода без условия существования. Основной конспект, раздел о силлогизме; саммари. & «Воссоздание…» (2010), PDF с. 1, печатная с. 88: автор приводит отсутствие вывода «Некоторые \(x\) суть \(y\)» из «Все \(x\) суть \(y\)» при двузначном толковании. Основной конспект отдельно проверяет схему с добавленной посылкой \(\exists t\,X(t)\). & Авторское утверждение плюс наш формальный вывод. & Подтверждено в указанной области: без посылки о существовании \(x\) частное суждение не следует; с такой посылкой классический вывод корректен. Это не свидетельствует, что двузначность сама по себе запрещает частное подчинение. \\
Н. П. Брусенцов считает, что закон исключённого третьего несправедливо приписан Аристотелю; проект не выдаёт это за установленную историю. Основной конспект, саммари. & «Трехзначная интерпретация…» (2003), PDF с. 1. Сопоставлено с Аристотелем, «Метафизика» \(\Gamma.7\), 1011b23–25: буквальный перевод по W. D. Ross приведён в основном конспекте. & Авторский исторический тезис плюс сопоставление с первичным текстом. & Цитата Б подтверждена. У Аристотеля есть текстовая формулировка отсутствия промежуточного между противоречащими; вопрос о названии закона как отдельной аксиомы оставлен отдельно. \\
Локальный PDF 2014 года неполон для второй публикации. Карта корпуса, таблица найденных расхождений. & Локальный PDF, страницы PDF 1–4: первая статья занимает печатные с. 239–241; вторая начинается на печатной с. 241 и обрывается в PDF с. 3; PDF с. 4 содержит только публикационную помету. Диапазоны с. 239–241 и 241–242 подтверждены библиографической записью \href{https://www.mathnet.ru/php/person.phtml?option_lang=rus\&personid=68115}{Math-Net}. & Сверка библиографии и фактического состава копии. & Исправлено: библиографические диапазоны сохранены, но копия второй статьи помечена как неполная; отсутствующее продолжение не пересказывается. \\
\bottomrule
\end{longtable}

\textbf{Общее правило применения списка.} Прямая цитата в первоисточнике подтверждает, что именно автор произнёс данный тезис; она не подтверждает автоматически его истинность. Формальный профиль (+-0+), таблица (+-++), кванторная запись конечного примера и высказывания о современном статусе модели помечаются отдельно как авторская формула, наш расчёт, наша формализация или исследовательская оценка. Строки с неполной пагинацией и семантикой считаются ограничениями источника, а не заполненными пробелами.

\subsection*{Формальная проверка критики импликации}

В таблице разведены объект и уровень каждого тезиса. Для проверки зафиксирована одна экстенсиональная реконструкция: непустой предметный домен \(D\), \(X,Y\subseteq D\), \(X'=D\setminus X\), \(Y'=D\setminus Y\); все кванторы пробегают один и тот же \(D\). Тогда материальная импликация на уровне предметов
\[
\forall u\,(X(u)\to Y(u))\;\equiv\;\neg\exists u\,(X(u)\land Y'(u))
\]
классически эквивалентна отсутствию контрпримерной клетки \(xy'\). Это выбранная экстенсиональная интерпретация записи автора, а не утверждение, что весь его формализм уже переведён.

{\small
\begin{longtable}{P{0.18\textwidth}P{0.15\textwidth}P{0.25\textwidth}P{0.25\textwidth}P{0.11\textwidth}}
\toprule
\textbf{Цитата / тезис и источник} & \textbf{Объект / формула} & \textbf{Область и предпосылки} & \textbf{Формальная проверка} & \textbf{Статус} \\
\midrule
\endfirsthead
\toprule
\textbf{Цитата / тезис и источник} & \textbf{Объект / формула} & \textbf{Область и предпосылки} & \textbf{Формальная проверка} & \textbf{Статус} \\
\midrule
\endhead
«Для несуществования \(xy'\)-вещей достаточно несуществования \(x\)-вещей либо \(y'\)-вещей»; строгая импликация Льюиса, 2005, PDF с. 1, печатная с. 140. & Отношение терминов через пустоту клетки: \(\neg\exists u(Xu\land Y'u)\). & Фиксированы один домен и классическое дополнение \(Y'=D\setminus Y\). Для предлагаемого автором отношения дополнительно нужны \(\exists uXu\) и \(\exists uY'u\). & Контрмодели: при \(D=\{a\}, X=\varnothing,Y=\varnothing\) клетка \(xy'\) пуста, но \(x\) также пуст; при \(X=Y=D\) пуста \(y'\). Следовательно, одна пустота \(xy'\) не обеспечивает требуемых условий сосуществования. После добавления условий существования получается именно более сильный критерий выбранной реконструкции. & Подтверждён локальный аргумент; он направлен против достаточности одной пустоты, а не против классической эквивалентности \(X\subseteq Y\leftrightarrow\neg\exists xy'\). \\
«Она равнозначна ему только в подчиненном ПСП универсуме Аристотеля УА»; отношение \(x\to y\) и необходимое \(x\Rightarrow y\), 2010, PDF с. 1–2, печатные с. 88–89. & Классовая импликация \(X\subseteq Y\) сопоставлена с отношением терминов; \(x\to y\equiv x'\lor y\equiv\neg(xy')\). & ПСП в данном экстенсиональном чтении требует \(X,X',Y,Y'\ne\varnothing\); следует зафиксировать, что штрих означает дополнение и что кванторы имеют общий домен. & По классической семантике \(X\subseteq Y\leftrightarrow\neg\exists(X\cap Y')\). ПСП выразим конечным набором предложений ФОЛ о непустоте классов. Например, \(D=\{a,b\}\), \(X=Y=\{a\}\), \(X'=Y'=\{b\}\) удовлетворяет ПСП и остаётся обычной бивалентной моделью. Это проверяет совместимость ПСП с ФОЛ, но не переводит авторские статусы и операции. & Различие формально заданных отношений; глобальная эквивалентность системы не доказана. \\
В 2006 году Н. П. Брусенцов утверждает, что двузначность не позволяет «погрузить» силлогистику и устранить парадоксы материальной импликации (PDF с. 1). & Тезис о выразительной и выводной способности логической системы в целом; не одна таблица \(P\to Q\). & Требуется указать класс моделей, значение статусов, допустимые непустые классы и отношение следования, относительно которого сравниваются системы. & Сам ПСП представим в классической ФОЛ, как показано выше; значит, он не противоречит бивалентности метаязыка. Однако из этого не следует, что все авторские операции и выводы сохраняются. Для этого нужен перевод с сохранением и отражением семантического следования либо выводимости. & Сильная общая оценка остаётся недоказанной этим местом; оснований диагностировать непонимание таблицы импликации нет. \\
Из «Все \(x\) суть \(y\)» не следует «Некоторые \(x\) суть \(y\)» при стандартном двузначном толковании; 2010, PDF с. 1, печатная с. 88. & Универсальная формула \(\forall u(Xu\to Yu)\) и частная \(\exists u(Xu\land Yu)\). & Домен непуст, но расширение \(X\) не обязано быть непустым. Если принята экзистенциальная предпосылка \(\exists uXu\), результат меняется. & При \(D=\{a\}\), \(X=\varnothing\), \(Y=\varnothing\), универсальная формула истинна, частная ложна. При добавлении \(\exists uXu\), из универсальной формулы классически следует частная. Причина — отсутствие экзистенциального импорта у выбранной семантики, а не бивалентность сама по себе. & Локальный пример верен для современной ФОЛ без экзистенциального импорта; спор касается допустимой семантики общего суждения и фоновых условий. \\
Табличная строка \(F\to T=T\) для материального условного часто называется парадоксальной; сопоставление с источниками 2005, 2006 и 2010 годов. & Пропозициональная формула целых высказываний \(P\to Q\equiv\neg P\lor Q\), отдельно от классового \(x\to y\). & При оценке \(P=F,Q=T\) материальное условное истинно по классической таблице; причинная, объяснительная или релевантная связь в формулу не входит. & Таблица устанавливает только истинность условного на этой оценке. Из неё не следует ни что \(P\) причиняет \(Q\), ни что \(P\) релевантно \(Q\). Ложная посылка как часть множества посылок в классической логике — ещё один вопрос и не тождественна истинности условного \(P\to Q\). & Формальная истина не опровергает философскую критику неадекватности для содержательной связи; смешивать эти объекты было бы ошибкой, но в просмотренных формулах автора такое смешение не установлено. \\
\bottomrule
\end{longtable}
}

В метаязыке классическое семантическое следование \(\Gamma\models Q\) означает: во всякой модели, удовлетворяющей всем посылкам из \(\Gamma\), истинно \(Q\). Поэтому противоречивое множество посылок семантически влечёт любое заключение из-за отсутствия моделей; это не оценка одной формулы \(P\to Q\). Отдельно, при оценке \(P=F,Q=T\), материальное условное истинно. Ни это табличное значение, ни отношение следования сами по себе не задают релевантную или объяснительную связь. В просмотренных местах Н. П. Брусенцов обсуждает прежде всего условия существования терминов и статус классового отношения; эти уровни не следует подменять друг другом.

\textbf{Вывод проверки.} На просмотренных страницах Н. П. Брусенцов явно задаёт материальную импликацию в терминах \(x'\lor y\) и отсутствия \(xy'\), а затем указывает, почему одна пустота клетки не даёт его отношения при отсутствии носителей терминов. Это текстовое основание считать его тезис содержательной критикой адекватности классической семантики принятой им задаче, а не свидетельством непонимания таблицы истинности. При этом широкие выводы 2006 года о двузначной логике не доказываются одной этой формулой: они требуют отдельного определения авторского отношения следования и сравнения систем относительно одного класса моделей.

\subsection*{Что именно означает «содержательное следование»}

Сравнение ниже исходит из определения самого автора. В главе 2005 года Н. П. Брусенцов задаёт необходимое следование \(x\Rightarrow y\) как отношение между терминами/классами, объясняет его через четыре клетки и профиль \(+-0+\) (PDF с. 1–3, печатные с. 140–142). В статье 2006 года выражение «содержательное следование» стоит в заглавии и относится к этому предложению; там же автор называет его трёхзначным и даёт алгебраическое представление (PDF с. 1). В 2010 году он пишет \((x\Rightarrow y)\equiv Axy\) и утверждает равнозначность материальной импликации с этим отношением только в универсуме УА (PDF с. 1–2, печатные с. 88–89). В рассмотренных публикациях формулировки меняются между «необходимым», «содержательным» и «аристотелевым необходимым» следованием. Поэтому «содержательное следование» — авторски засвидетельствованное название, но не устоявшийся синоним любого логического следования. Для формальной проверки ниже употребляется более узкое обозначение: \emph{предлагаемое Н. П. Брусенцовым отношение необходимого следования между терминами} \(x\Rightarrow y\); отдельные публикации не отождествляются без проверки.

\textbf{Различение по типу объекта.} У Челпанова в скане пятого издания 1909 года силлогизм определяется как форма умозаключения, где из двух суждений необходимо вытекает третье; при этом он отдельно различает материю (термины) и форму силлогизма (печатные с. 75–78). Для условного силлогизма формула «если A есть B, то C есть D» состоит из основания и следствия (с. 97–98). Доказательство у Челпанова — выведение суждения из других суждений, признанных истинными и очевидными (с. 144). Значит, необходимость вывода, связь частей условного суждения и доказательство с истинными основаниями — разные предметы. Пагинация здесь печатная; изображения кадров 87–88, 109–110, 119–120 и 156 просмотрены в скане НБУ. Это ограниченный набор глав, а не полная транскрипция учебника.

{\small
\setlength{\tabcolsep}{3pt}
\begin{longtable}{P{0.18\textwidth}P{0.28\textwidth}P{0.18\textwidth}P{0.23\textwidth}}
\toprule
\textbf{Термин / тип объекта} & \textbf{Определение / критерий} & \textbf{Источник и место} & \textbf{Совпадение с \(x\Rightarrow y\)} \\
\midrule
\endfirsthead
\toprule
\textbf{Термин / тип объекта} & \textbf{Определение / критерий} & \textbf{Источник и место} & \textbf{Совпадение с \(x\Rightarrow y\)} \\
\midrule
\endhead
«Содержательное следование» Н. П. Брусенцова\newline\emph{Авторская характеристика отношения между терминами/классами} & Необходимое отношение \(x\Rightarrow y\); в конкретных работах задаётся своими клеточными статусами, алгеброй и предпосылками. & 2005, PDF с. 1–3; 2006, PDF с. 1; 2010, PDF с. 1–2. & Базовое авторское отношение. Само название не задаёт семантику вне указанной работы. \\
Четырёхпозиционный профиль Н. П. Брусенцова \(+-0+\)\newline\emph{Статусная запись отношения терминов} & В порядке \(xy,xy',x'y,x'y'\): необходимое включение, исключение, привходящий статус, включение. & «Реанимация аристотелевой силлогистики» (2005), PDF с. 1–3, печатные с. 140–142. & Исходный объект сравнения: статусы сочетаний терминов, а не маркеры занятости четырёх областей или значения целых пропозиций. \\
Малая диаграмма Льюиса Кэрролла\newline\emph{Разбиение классов вещей по двум признакам} & Область — класс предметов с определённым сочетанием признаков; красная фишка означает «занята», серая — «пуста», пограничная фишка — занятость одной из двух смежных областей при неизвестном выборе. & Lewis Carroll, \emph{The Game of Logic} (1886), печатные с. 4–5, 9; изображения страниц проверены. & Геометрически те же четыре сочетания признаков; значения маркеров отличаются от статусов \(+-0+\). Оригинал не приписывает каждой клетке профиль необходимых отношений. \\
Диаграмма Кэрролла в интерпретации Н. П. Брусенцова\newline\emph{Авторская модификация диаграммного кодирования} & В подписанном Н. П. Брусенцовым фрагменте 2011 года 1 обозначает существование, 0 — несуществование/несовместимость; автор читает незанятую клетку \(x'y\) в УА как сохранение третьего статуса. & «Алгебраизация безупречной логики Аристотеля…» (2011), локальный фрагмент PDF с. 1–3, печатные с. 3–5; изображения просмотрены. Книга целиком — 12 с.; локальная копия фрагментарна. & Совпадение с кэрролловой схемой заявлено как конкретная авторская конструкция; совпадение четырёх областей само по себе не доказывает эквивалентность всех операций. \\
Условное суждение у Челпанова\newline\emph{Суждение с основанием и следствием} & «Если A есть B, то C есть D»; modus ponens утверждает основание и заключает следствие. & Челпанов, 5-е изд. (1909), печатные с. 97–98; скан НБУ. & Совпадения нет: целое условное суждение и его части не являются отношением между терминами \(x,y\). \\
Категорический силлогизм\newline\emph{Дедуктивный вывод из суждений} & При заданной форме заключение необходимо следует из посылок; правильность формы отделена от истинности посылок. & Челпанов, 1909, печатные с. 75–78; Аристотель, \emph{Первая аналитика} I.1, 24b18–22. & Частичное сходство, но не совпадение: \(Axy\) у автора — суждение/отношение терминов, используемое в силлогизме, а не сам двухпосылочный вывод. \\
Демонстрация / доказательство\newline\emph{Доказательный и познавательный процесс} & У Челпанова — выведение из признанных истинными и очевидными суждений. У Аристотеля научная демонстрация требует истинных, первичных, непосредственных, предшествующих и причиняющих посылок. & Челпанов, 1909, с. 144; Аристотель, \emph{Вторая аналитика} I.2, 71b9–72b4. & Совпадения нет: отношение \(x\Rightarrow y\) само не задаёт ни доказательство, ни статус знания, ни причину заключения. \\
Материальная импликация\newline\emph{Связка для целых высказываний; вторая формула — классовая} & Для целых высказываний \(P\to Q\equiv\neg P\lor Q\). В классовом чтении \(x\to y\equiv x'\lor y\equiv\neg(xy')\). & Н. П. Брусенцов, 2010, PDF с. 1–2; современное описание: SEP, «Relevance Logic», § 1. & Условное совпадение по тезису автора только в УА; вне этой предпосылки общая эквивалентность не установлена. \\
Семантическое следование \(\Gamma\models\varphi\)\newline\emph{Метаязыковое отношение посылок и заключения} & Во всякой допустимой модели, удовлетворяющей всем формулам \(\Gamma\), истинно \(\varphi\); критерий — отсутствие контрмодели. & SEP, «Logical Consequence», §§ 3.1–3.2. & Не определено как авторское отношение: тексты задают отношение терминов и процедуру для силлогистических посылок, но не общий класс \(\Gamma\) и критерий для произвольных \(\varphi\). \\
Строгая импликация К. И. Льюиса\newline\emph{Модальное отношение между целыми высказываниями} & \(P\) строго влечёт \(Q\), если невозможно совместное положение \(P\land\neg Q\); современная запись — \(\neg\Diamond(P\land\neg Q)\). & C. I. Lewis, \emph{A Survey of Symbolic Logic} (1918), гл. V, печатные с. 293, 324; SEP, «Clarence Irving Lewis». & Совпадения с диаграммой классов нет: это отношение между пропозициями через модальное понятие невозможности. В корпусе Н. П. Брусенцова критика строгой импликации отделена от тезиса о материальной импликации. \\
Трёхзначная логика Яна Лукасевича (1920)\newline\emph{Матрица значений пропозиций} & В первичной статье: 0 — ложь, 1 — истина, 2 — «возможность»; современные изложения нормируют третье значение как \(\tfrac12\). & «O logice trójwartościowej», \emph{Ruch Filozoficzny}, т. 5, № 9 (1920), печатные с. 170–171, скан просмотрен; SEP, § 5.1. & Совпадения с профилем ячеек нет по типу объекта: это значения целых суждений и их связок, а не занятость классовых областей и не значение нуля в профиле 2005 года. \\
Релевантная импликация / логика релевантности\newline\emph{Семейство связок и исчислений для связи посылок и вывода} & В частности, принцип общих пропозициональных переменных — необходимое, но недостаточное условие; некоторые системы явно отслеживают использование предпосылок в доказательстве. & SEP, «Relevance Logic», §§ 1, 4. & Совпадения нет по рассмотренным критериям: основа у Н. П. Брусенцова — статусы терминальных клеток и ПСП; общие переменные или использование предпосылки здесь не заданы. \\
Выводимость \(\Gamma\vdash\varphi\)\newline\emph{Отношение доказательства в фиксированном исчислении} & Существует конечное доказательство \(\varphi\) из \(\Gamma\) по правилам системы. & SEP, «Logical Consequence», § 3.2; глава 2005 года, PDF с. 5–6. & Частично процедурно: автор задаёт алгоритм пересечения шкал и устранения среднего термина для модусов; это не полная эквивалентность произвольной выводимости. \\
Фактическое использование каждой посылки\newline\emph{Свойство доказательства / зависимости от предпосылок} & Посылка используется, если её допущение действительно участвует в данном доказательстве; это не то же, что истинность или семантическая незаменимость посылки. & SEP, «Relevance Logic», § 4; авторская процедура 2005 года, PDF с. 5–6. & Не определено: элиминация среднего термина описана, но общего критерия «каждая посылка использована» не найдено. \\
Причинная / объяснительная связь\newline\emph{Отношение объяснения между основаниями и фактом} & В аристотелевской демонстрации посылки не просто влекут заключение, но должны быть его причинами. & Аристотель, \emph{Вторая аналитика} I.2, 71b9–72b4; Н. П. Брусенцов, 2008, PDF с. 1. & Совпадения нет по имеющимся определениям: объяснение через сущность дано как философское основание автора, но \(x\Rightarrow y\) не определено как причинная связь. \\
\bottomrule
\end{longtable}
}

\textbf{Итог терминологической проверки.} В проверенных публикациях «содержательное следование» — собственное название и характеристика предложения Н. П. Брусенцова, но не синоним материальной или строгой импликации, семантического следования, выводимости, релевантной связи либо причинного объяснения. Частичное процедурное сходство есть у классического силлогизма и у требования использовать посылки в релевантном доказательстве; это не тождество формализмов. Оригинальность можно оценивать только после пооперационного сравнения; имеющийся выборочный просмотр корпуса не устанавливает ни приоритета, ни полного отсутствия предшественников.

\section{Подтверждённые положения}

\subsection{Термины не являются пропозициональными буквами}

В статьях о силлогистике \(x\) и \(y\) обозначают признаки/термины и связанные с ними классы предметов; \(xy,xy',x'y,x'y'\) — четыре сочетания признаков. Поэтому знак \(x\Rightarrow y\) у Н. П. Брусенцова относится к предлагаемому отношению между терминами. Его нельзя без объяснения заменить формулой \(P\to Q\), где \(P,Q\) — целые высказывания.

\subsection{Схема четырёх клеток и её область применения}

В главе 2005 года профиль \(+-0+\) дан в порядке \(xy,xy',x'y,x'y'\). Знаки задают необходимое включение, необходимое исключение и привходящий статус. Для позиции \(x'y\) автор пишет, что подкласс не включён и не исключён с необходимостью: его статус зависит от иных обстоятельств.

В главе 2005 года Н. П. Брусенцов сам объясняет экстенсиональный смысл своей нормальной формы: при сосуществовании противоположностей пустота (xy') в непустом классе (x) влечёт непустоту (xy), а в непустом классе (y') — непустоту (x'y'). Вектор (+-0+) затем кодирует необходимые включение/исключение этих сочетаний и привходящий статус (x'y). Поэтому авторское толкование статусов через непустоту классов опирается на источник. Однако экстенсиональная формула с кванторами, приведённая в конспекте, требует непустоты xy и x'y', пустоты xy' и ничего не говорит о фактической пустоте/непустоте x'y; это лишь экстенсиональная тень полной статусной формы. Сам символ (+) не является квантором, а применение авторского объяснения к конечной предметной области — наша запись, не буквальная формула Н. П. Брусенцова.

Если отдельно прочитать материальную импликацию \(X(t)\to Y(t)\) по четырём комбинациям принадлежности, получится профиль \(+-++\): ложна только \(xy'\). Это наш точечный пересчёт обычной импликации для предикатного случая, не профиль, который Н. П. Брусенцов вводит в источнике. Его профиль \(+-0+\) дан для классовых выражений и задаёт привходящий статус комбинации \(x'y\), тогда как экстенсиональная формула, свободно оставляющая эту клетку, не передаёт данный статус. Это различие способов оценки предметных сочетаний, но оно не доказывает, что семантическое следование \(\Gamma\models Q\) или правила вывода неверны.

\subsection{В источниках 2003 и 2005 годов разные кодировки}

В статье 2003 года функция \(A(x,y)\) использует \(1,0,\sigma\): необходимую присущность, антиприсущность и привходящий статус. Автор указывает \(0<\sigma<1\), но не задаёт конкретное числовое значение и не уточняет, допускаются ли разные числовые реализации этого статуса. В главе 2005 года \(\sigma\) прямо названо третьим значением истинности обобщённого выражения о классах, а знаковая алгебра кодирует необходимые включение и исключение как \(+1\) и \(-1\), привходящий статус — как \(0\). Поэтому:

\begin{itemize}
\item третье значение у Н. П. Брусенцова относится к его выражениям о классах;
\item нулевой знак в третьей позиции \(+-0+\) отмечает привходящий статус конкретной комбинации \(x'y\);
\item ни одно из этих утверждений не означает третье значение таблицы \(P\to Q\);
\item число \(0\) в \(A(x,y)\) за 2003 год и число \(0\) в знаковой шкале 2005 года имеют разные назначенные смыслы.
\end{itemize}
\subsection{Сосуществование противоположностей — дополнительная предпосылка}

В статье 2003 года и последующих текстах Н. П. Брусенцов требует, чтобы первичный термин и его противоположность имели экземпляры в универсуме. В статье 2003 года он прямо отмечает, что Аристотель не сформулировал это условие явно. Его следует описывать как дополнительный постулат реконструкции Н. П. Брусенцова, а не как установленную аксиому Аристотеля.

Условие существенно для выводов: общая фраза о всех \(x\) может оказаться истинной из-за пустоты класса \(x\). Чтобы перейти от общего утверждения к частному, существование \(x\)-предметов необходимо указать отдельно.

\subsection{Реализация операций не устанавливает их исторический статус}

В публикациях о троичной компьютеризации и конструктном кодировании описаны шкалы, операции и получение силлогистических выводов. Это показывает, что заданную систему можно алгоритмически представить и исполнять. Само по себе это не доказывает, что её семантика исторически принадлежит Аристотелю или должна заменить все современные логики.

\section{Выразимость в классической ФОЛ и предварительная оценка новизны}

Слово «сведение» допускает несколько разных критериев: кодирование конечных статусов как данных, перевод формул с квантификацией по ситуациям, сохранение таблиц операций или эквивалентность отношения следования. Успех первого или второго критерия не доказывает автоматически третий. Классическая логика первого порядка может описывать в одном языке ситуации и предметы, но такое описание ещё должно быть доказано семантически корректным относительно конкретного исходного исчисления.

\begin{longtable}{P{0.22\textwidth}P{0.36\textwidth}P{0.30\textwidth}}
\toprule
\textbf{Проверяемый тезис} & \textbf{Сверенное свидетельство} & \textbf{Что из него следует} \\
\midrule
\endfirsthead
\toprule
\textbf{Проверяемый тезис} & \textbf{Сверенное свидетельство} & \textbf{Что из него следует} \\
\midrule
\endhead
Троичное различение вообще нельзя представить двузначными средствами. & В «Математической теории силлогистики» Н. П. Брусенцов пишет, что разбиение на «присущее», «неприсущее» и «неопределено» можно смоделировать двузначной логикой в два шага: сначала выделить неопределённый класс, затем разбить остаток на два класса (PDF с. 4). Изображение страницы просмотрено. & Универсальный запрет на всякое двузначное кодирование противоречит этому раннему замечанию самого автора. Оно касается кодирования разбиения, а не всей поздней алгебры следования. \\
Трёхзначная логика как таковая впервые появляется в корпусе Н. П. Брусенцова. & Трёхзначное исчисление высказываний Яна Лукасевича опубликовано в 1920 году; позднейшие изложения приводят формальную матрицу и её пропозициональную семантику. В работах Н. П. Брусенцова 2003/2005 годов три статуса относятся к отношениям терминов/классовых ячеек, а не автоматически к значениям целых высказываний. & Само наличие трёх значений не является новым. Возможная оригинальность относится к типу объектов, их четырёхпозиционному профилю, операциям и силлогистическому кодированию; прямое тождество двух систем не утверждается. \\
Как посылки производят силлогистическое заключение. & В главе 2005 года две посылки задаются трёхтерминными шкалами; автор пересекает эти шкалы и элиминирует средний термин, чтобы получить заключение, если оно существует (PDF с. 5). Там же разобраны Barbara, Bamalip и другие модусы (PDF с. 5–6). В статье 2003 года сказано, что заключение достигается преобразованием конъюнкции посылок (PDF с. 6). Изображения страниц 2005 с. 5–6 и 2003 с. 6 просмотрены. & В корпусе есть конструктивная, композиционная процедура вывода; вопрос не сводится к выразимости отдельного профиля. Это ещё не тождественно условию семантической незаменимости каждой посылки и не доказывает эквивалентность процедуры классическому семантическому следованию. Для такой эквивалентности нужен точный перевод операций и доказательство сохранения выводов. \\
Классическая булева алгебра не умеет выразить смысл поздних статусных операций. & В «Конструктной компьютеризации силлогистики» (2007, печатная с. 10; PDF с. 1) Н. П. Брусенцов и Ю. С. Владимирова ограничивают тезис обсуждением булевой алгебры, конъюнктивных/дизъюнктивных форм и пропуска несущественного члена в ДНФ/КНФ. Соответствующая страница просмотрена. & Это аргумент об адекватной композиционной записи внутри обсуждаемой алгебры. Он сам по себе не является теоремой о невозможности кодирования в многосортной ФОЛ с предикатами статусов или переменной ситуации. \\
Как именно двузначная ДНФ искажает несущественную клетку? & В статье «Аристотелево необходимое следование» (2014, PDF с. 2; печатная с. 240) Н. П. Брусенцов пишет, что клетка в диаграмме Кэрролла остаётся незаполненной и может быть как пустой, так и занятой; его способ полной записи трёхзначного отношения опустошает соответствующую единичную клетку материальной импликации. Изображение первичной страницы просмотрено; текстовый слой этой копии имеет ошибки кодировки. & Это уточняет авторскую критику конкретной двузначной ДНФ и показывает совпадение проблемы с потерей статуса в нашей экстенсиональной тени. Это не доказательство невозможности кодирования в классической ФОЛ с переменной ситуации. Вторая статья 2014 года, начинающаяся на печатной с. 241, представлена в локальной копии неполно. \\
Ситуационные статусы можно записать формулами классической ФОЛ. & В основном конспекте построены формулы непустоты четырёх клеток, устойчивой непустоты/пустоты по доступным ситуациям и изменения статуса. Условия на непустую рамку делают три статуса исчерпывающими и несовместимыми; для смешанного статуса предъявлены разные оценки клетки в двух доступных ситуациях. & Доказана выразимость именно выбранной интерпретации «необходимо / исключено / меняется по ситуациям». Это реконструкция, а не приписанное источнику определение \(\sigma\); авторское отношение из статьи 2003 года и его операции этой выразимостью не покрыты. \\
Такой перевод уже сводит всё следование Н. П. Брусенцова к классическому. & В переводе с по-ситуационной конъюнкцией найден контрпример сохранению минимума: два смешанных статуса могут быть взаимно исключены и их конъюнкция окажется всегда ложной. Более того, одна и та же пара входных статусов \((\sigma,\sigma)\) может быть получена из смешанной оценки \(p\) и \(p\) или из \(p\) и \(\neg p\), но поточечные конъюнкция/дизъюнкция в этих случаях дают разные статусы. Поэтому детерминированная таблица только на трёх сжатых статусах не воспроизводит обе операции в произвольной ситуативной модели. & Это ограничение выбранного поточечного перевода, не опровержение авторской таблицы. Для полной редукции нужны композиционный перевод всех операций и правил, полный список свойств и доказательства их сохранения и отражения. Такого результата ревизия не устанавливает. \\
\bottomrule
\end{longtable}

Стандартный перевод реляционной модальной семантики в формулы первого порядка с переменной состояния известен в модальной логике (Blackburn и van Benthem, 2007, раздел 2.2). Кроме того, аристотелевскую немодальную силлогистику описывают как фрагмент первого порядка (SEP, «Ancient Logic», раздел 2.4). Эти результаты делают сам метод ситуационного кодирования обычным формальным инструментом, но не решают вопрос о корректности именно модели Н. П. Брусенцова.


\subsection{Четыре разных критерия участия посылок}

Для проверки «содержательной связи» нужно заранее указать, какого рода связь ищем. Сравним четыре аналитических критерия на Barbara из главы 2005 года, \(A_{yz},A_{xy}\Rightarrow A_{xz}\).

\begin{longtable}{P{.18\linewidth}P{.18\linewidth}P{.32\linewidth}P{.24\linewidth}}
\toprule
\textbf{Критерий} & \textbf{Объект} & \textbf{Проба} & \textbf{Связь с источником} \\
\midrule
\endfirsthead
\toprule
\textbf{Критерий} & \textbf{Объект} & \textbf{Проба} & \textbf{Связь с источником} \\
\midrule
\endhead
Семантическая незаменимость относительно \(T\) & Следование \(T\cup\Gamma\models C\), где \(T\) — условия сосуществования терминов и противоположностей. & В нашей экстенсиональной ФОЛ-реконструкции обе посылки необходимы: если удалить \(A_{xy}\), взять \(X=\{a\}\), \(Y=Z=\{b\}\); если удалить \(A_{yz}\), взять \(X=Y=\{a\}\), \(Z=\{b\}\). Обе контрмодели удовлетворяют \(T\); оставшаяся посылка верна, а \(A_{xz}\) неверно. & Это проверка незаменимости по выбранному переводу и фоновым условиям, не авторский критерий использования посылок и не структурное свойство доказательства. \\
Использование в доказательстве & Строение вывода \(\Gamma\vdash C\). & В процедуре 2005 года обе шкалы входят в пересечение; затем удаляется общий \(y\). Это показывает использование обеих посылок именно этим алгоритмом. & Автор описывает вычислительный шаг на с. 5; отдельного исчисления с недисчардженными допущениями в просмотренном месте нет. \\
Связность через средний термин & Структура терминов/посылок. & В Barbara обе посылки соединены через \(y\), который затем элиминируется. Но наличие такого пути само по себе не гарантирует необходимость посылок: в \(A_{xy},A_{yz},A_{xz}\vdash A_{xz}\) прямая посылка \(A_{xz}\) делает путь через \(y\) избыточным. & Элиминация \(y\) — прямое описание автора. Более общий критерий релевантности этим ещё не задан. \\
Объяснение через сущность & Объяснительное отношение, а не только выводимость. & Вычисление Barbara даёт формальную связь \(x\Rightarrow z\), но одно оно не устанавливает, что \(z\) объясняет \(x\) или входит в его сущность. & В работе 2008 года автор связывает содержательность с тем, что качество \(y\), присущее \(x\)-вещам, необходимо содержится в их сущности (PDF с. 1). Это авторское философское основание, а не формальный критерий использования посылок. \\
\bottomrule
\end{longtable}

Таким образом, пример Barbara подтверждает, что описанная Н. П. Брусенцовым процедура принимает обе шкалы, пересекает их и устраняет средний термин. Мы проверили два вычислительных примера по своей реконструкции этой процедуры, но соответствие её общему формальному правилу пока не установлено. Семантическая незаменимость в приведённых моделях относительно \(T\) — отдельный результат нашей реконструкции; она не равна использованию посылки в доказательстве: например, в классической логике \(\{P,\neg P\}\models Q\) может быть минимальным по включению набором посылок, хотя такая связь с \(Q\) нерелевантна. Следование \(\models\) монотонно; если авторский критерий участия посылок исключает сохранение вывода при добавлении посторонней посылки, он не совпадёт с классическим семантическим следованием. Однако такой критерий сначала нужно установить в первоисточнике. Объяснение через сущность — отдельное авторское основание. Ни один из этих фактов сам по себе не задаёт целевое авторское отношение следования или причинное объяснение.

\textbf{Предварительный вывод о новизне.} Проверенные материалы не поддерживают сильный тезис о принципиальной невыразимости трёх статусов в классической ФОЛ вообще: для функции отношений 2003 года её качественные статусы выражаются формулами о непустоте двух клеток при авторском постулате, а сам Н. П. Брусенцов в 1971 году допускает двузначное кодирование трёхчастного разбиения. В 2014 году он обсуждает более конкретную недостаточность двузначной ДНФ для представления несущественной клетки. Ни один источник не доказывает эквивалентности позднейших операций и правил классической логике первого порядка; и выразимость отдельных статусов сама по себе этого не доказывает. Проверенная часть корпуса показывает предложенную автором алгебраическую реконструкцию и процедуру получения ряда силлогистических заключений. Но оснований объявить её ни «новой логикой», ни простой «переупаковкой» пока нет: исторический приоритет и теоретическое новшество требуют пооперационного сравнения с конкретными системами по общему критерию. Вывод ограничен просмотренными публикациями и страницами, а не полным массивом предшествующих работ.

\section{Найденные исправления и открытые расхождения}

\begin{longtable}{P{0.260\textwidth}P{0.300\textwidth}P{0.360\textwidth}}
\toprule
\textbf{Тема} & \textbf{Результат ревизии} & \textbf{Действие} \\
\midrule
\endfirsthead
\toprule
\textbf{Тема} & \textbf{Результат ревизии} & \textbf{Действие} \\
\midrule
\endhead
\(\sigma\) и «третье значение» & Автор прямо употребляет это выражение для классового выражения; в 2005 году привходящий статус кодируется нулём. & Уточнён основной конспект и саммари. \\
Функция статусов статьи 2003 года & PDF-копия, с. 1–4: при непустоте всех четырёх противоположных терминов значения 1, 0 и \(\sigma\) выражаются соответственно пустотой \(xy'\), пустотой \(xy\) и непустотой обеих клеток. Вне этого условия статус не задан; проверены изображения PDF-с. 1–4, пагинация электронная, копия производная от архивного DOC. & Добавлена классическая формализация для фиксированной совокупности и оговорено, что она не доказывает сохранения операций или процедуры вывода. \\
Ноль в позиции \(x'y\) & Это привходящий статус именно этой комбинации, которая может быть включена или исключена; на конечную группу объектов это не переносится без отдельного толкования. & Уточнены карта корпуса и пояснение к примеру. \\
Материальная импликация & Вектор \(+-++\) получается при нашем пересчёте обычной импликации по четырём сочетаниям терминов; профиль \(+-0+\), предложенный Н. П. Брусенцовым, дан в алгебре классовых выражений. & Сравнение помечено как реконструкция, а не прямая цитата или таблица истинности, предложенная Н. П. Брусенцовым. \\
Число правильных модусов & Статья 2003 года считает \(4\times8^3=2048\) комбинаций и сообщает 192 (PDF с. 6), английская аннотация той же производной копии — 190 (PDF с. 10). Статья 2007 года сообщает 128 при кодировании четырёхтритными конструктами и склеивании \(\oplus\) (PDF с. 2); статья 2008 года также сообщает 128 (PDF с. 6). & Общий список различаемых суждений и единый критерий правильности не восстановлены; числа оставлены как расхождение источников, не как пересчёт. \\
Статьи 2014 года & Публикационная запись указывает первую статью на с. 239–241, вторую на с. 241–242. Локальный PDF показывает первую статью и только начало второй на с. 241; её продолжения на с. 242 в копии нет. & В карте сохранены библиографические диапазоны и отдельно отмечена неполнота локальной копии. \\
Ссылки между публикациями и число 128 & Статья ММРО-13 (2007) ссылается на публикацию 2003 года; в списке литературы статьи 2008 года ссылки на ММРО-13 не найдено. & Неоднозначная фраза в карте заменена точным описанием найденных ссылок. \\
\bottomrule
\end{longtable}

Числа 128, 190 и 192 нельзя объединить только вычислением, пока не восстановлено, одинаково ли авторы считали фигуры, формы, перестановки, пустые классы и условия правильности. Lean может пересчитать число для одной формальной классификации, но не выведет из текстов, какую именно классификацию подразумевал каждый автор.

\section{Что статья 2010 года определяет в своих формулах}

В статье «Воссоздание аристотелевой безукоризненной силлогистики» (2010, PDF, с. 1; печатная с. 88) напечатано:

\( (x\Rightarrow y) \equiv Axy \equiv I_{xy} I'_{xy'} I_{x'y'}. \)

Автор называет \(x\Rightarrow y\) необходимым следованием и связывает \(Axy\) с общеутвердительным суждением «Все \(x\) суть \(y\)». На той же PDF-с. 1 он поясняет исключение сочетания \(xy'\) через формулу \(x\to y \equiv (x'\lor y) \equiv \neg(xy') \equiv I'_{xy'}\). Здесь стрелка \(x\to y\) относится к материальной импликации в алгебре терминов, а не автоматически к пропозициональной формуле \(P\to Q\). Текст даёт контекстное прочтение для \(I'_{xy'}\), но не полный словарь обозначений.

\textbf{Статья локально не определяет каждый символ и операцию.} Парная статья «Трехзначность отношения следования» (2010, печатные с. 86–87) поясняет \(I_{xy}\) через возможность отношения «Некоторые \(x\) суть \(y\)»; длинная статья даёт для \(I'_{xy'}\) чтение через исключение контрпримера. Это межтекстовое уточнение, но не полное определение каждого множителя: в статье 2003 года сходное обозначение \(Ixy\) входит в классовое суждение и связано с \(Vxy\) при условии принятого универсума. Поэтому сведение всех множителей к простым формам существования остаётся реконструкцией. Правило соседнего соединения множителей и точная семантика каждого знака \(\equiv\) в статье также не задаются.

\textbf{Сомнение в последнем штрихе разрешено визуально.} В формуле стоит \(I_{x'y'}\), а не \(I'_{x'y'}\): штрихи относятся к обоим индексам, но не к оператору \(I\). Это подтверждается изображением PDF-с. 1, диаграммой PDF-с. 3 и параллельными формулами в работах 2005, 2006 и 2011 годов. Следовательно, расхождение возникло из-за малого размера знака при извлечении/прочтении, а не из-за разных вариантов формулы в источниках.

Проверена также формула на PDF-с. 2 (печатная с. 89), которую прежняя версия руководства не разбирала:
\[
Axy\,I'_{x'y}
\equiv AxyAyx
\equiv I_{xy}I'_{xy'}I'_{x'y}I_{x'y'}
\equiv (x\equiv y).
\]
Текст связывает \(I'_{x'y}\) с исключением случая «\(y\), но не \(x\)» и поясняет, что одновременные \(Axy\) и \(Ayx\) дают тождество терминов. Это дополняет первую формулу: для \(Axy\) само по себе \(x'y\) остаётся привходящим, а обратное \(Ayx\) требует его исключить. Транскрипция знаков проверена по изображению; формальное правило для соседнего соединения факторов по-прежнему не задано.

В публикации «Трехзначная диалектическая логика» (2001) автор называет \(V\) интегральной дизъюнкцией — специальным оператором своей модели, не обычным \(\lor\) — и отдельно записывает условие для терминов и их противоположностей:

\( V_x V_{x'} V_y V_{y'} \equiv 1. \)

Это условие относится к существованию каждого первичного термина и его противоположности. Оно не совпадает с клеточной записью

\( Axy \equiv V'_{xy'} V_x V_{x'} V_y V_{y'} \equiv V_{xy} V'_{xy'} V_{x'y'}, \)

которую Н. П. Брусенцов поясняет через существование сочетаний \(xy\) и \(x'y'\), исключение \(xy'\), а также привходящий статус \(x'y\). Следовательно, экстенсиональное прочтение профиля \(+-0+\) в модели 2005 года опирается на прямое объяснение Н. П. Брусенцова; при этом сам символ \(+\) не является квантором существования, а применение его условий к конечному списку — предметная формализация. Публикация 2010 года не утверждает, что \(I\) и \(V\) тождественны; межтекстовая параллель делает такое сопоставление правдоподобным, но не устанавливает тождество, поэтому не следует автоматически переносить все условия существования на формулу 2010 года. Подробное чтение обозначений и наглядный пример приведены в \href{../docs/Как-читать-формулы-Брусенцова.tex}{руководстве по формулам}.

\section{Источниковые ограничения}

\begin{itemize}
\item Статья 2003 года доступна в архиве музея как исходный DOC; локальный PDF получен конвертацией. Это не факсимиле журнального издания, и электронная пагинация отличается от печатной.
\item В PDF 2003 года текстовые извлечения могут терять штрихи и логические символы; диаграммная страница содержит повреждённые/наложенные подписи. Формулы нужно сверять по изображениям, а не только по OCR.
\item В PDF 1982 года нет надёжного текстового слоя. Рецензент проверил, что файл является сканом и читаем визуально, но не транскрибировал каждую формулу всех страниц.
\item Полная книга \emph{Реставрация логики} в проверенных открытых каталогах не найдена. Доступна отдельная глава, а не весь том.
\item Библиографический указатель охватывает 207 нумерованных записей списка МГУ на 20.11.2026. Это полный индекс относительно данного списка; он не гарантирует полноту всех публикаций, изобретений и изданий.
\end{itemize}
\section{Общий вывод ревизии}

Концептуальное ядро корпуса теперь описано точнее: Н. П. Брусенцов предлагает алгебраическую модель отношений между терминами и классами, вводит условие существования противоположных классов и использует привходящий статус. Это содержательное предложение можно проверять на внутреннюю непротиворечивость и корректность формальных выводов. Его оценка как решения проблемы содержательного следования, а также историческое прочтение Аристотеля остаются отдельными философскими и историческими задачами.

Для дальнейшей проверки полезно сначала закрепить одно полное определение системы и пересчитать её правила на малых моделях; только после этого имеет смысл строить более масштабную проверку в Lean.

\section{Библиография ревизии}

\subsection{Первичные работы Н. П. Брусенцова}

\begin{itemize}
\item Н. П. Брусенцов. \href{https://computer-museum.ru/histussr/trilog0.htm}{«Трехзначная интерпретация силлогистики Аристотеля» (2003)}. \href{../sources/pdfs/Брусенцов_Трехзначная-интерпретация-силлогистики-Аристотеля_2003_версия-музея.pdf}{Производная PDF-копия в корпусе}.
\item Н. П. Брусенцов. \href{https://www.computer-museum.ru/books/setun/brusencov_03.pdf}{«Математическая теория силлогистики» (1971), \emph{Вычислительная техника и вопросы кибернетики}, вып. 8, с. 154–176}. \href{../sources/pdfs/Брусенцов_Математическая-теория-силлогистики_1971.pdf}{PDF в корпусе}.
\item Н. П. Брусенцов. \href{https://ternarycomp.cs.msu.ru/Papers/Rea.pdf}{«Реанимация аристотелевой силлогистики» (2005)}. \href{../sources/pdfs/Брусенцов_Реанимация-аристотелевой-силлогистики_2005-глава-из-Реставрации-логики.pdf}{PDF в корпусе}.
\item Н. П. Брусенцов, Ю. С. Владимирова. \href{https://ternarycomp.cs.msu.ru/Papers/MMRO12.pdf}{«Троичная компьютеризация логики» (2005)}. \href{../sources/pdfs/Брусенцов_Владимирова_Троичная-компьютеризация-логики_2005_ММРО-12.pdf}{PDF в корпусе}.
\item Н. П. Брусенцов, Ю. С. Владимирова. \href{https://ternarycomp.cs.msu.ru/Papers/MMRO13.pdf}{«Конструктная компьютеризация силлогистики» (ММРО-13, 2007)}. \href{../sources/pdfs/Брусенцов_Владимирова_Конструктная-компьютеризация-силлогистики_2007_ММРО-13.pdf}{PDF в корпусе}.
\item Н. П. Брусенцов. \href{https://ternarycomp.cs.msu.ru/Papers/Alrek.pdf}{«Алгебраическая реконструкция силлогистики» (2008)}. \href{../sources/pdfs/Брусенцов_Алгебраическая-реконструкция-силлогистики_2008.pdf}{PDF в корпусе}.
\item Н. П. Брусенцов. «Трехзначность отношения следования» (2010), печатные с. 86–87. \href{https://ternarycomp.cs.msu.ru/Papers/brusentsov_sborn2010_1.pdf}{Источник на сайте МГУ}; \href{../sources/pdfs/Брусенцов_Трехзначность-отношения-следования_2010.pdf}{локальный PDF}.
\item Н. П. Брусенцов. \href{https://ternarycomp.cs.msu.ru/Papers/brusentsov_sborn2010_2.pdf}{«Воссоздание аристотелевой безукоризненной силлогистики» (2010), печатные с. 88–91}. \href{../sources/pdfs/Брусенцов_Воссоздание-аристотелевой-силлогистики_2010.pdf}{Локальная копия в корпусе}.
\item Н. П. Брусенцов. \href{https://ternarycomp.cs.msu.ru/Papers/Brusen_IMI14_Neobh_sled.pdf}{«Аристотелево необходимое следование» (2014), с. 239–241; «Трехзначное обобщение алгебры логики. Преодоление несовершенности ДНФ трехзначным обобщением логики» (2014), с. 241–242}. \href{../sources/pdfs/Брусенцов_Аристотелево-необходимое-следование-и-ДНФ_2014.pdf}{Общий PDF в корпусе}; вторая статья представлена только началом, продолжение с. 242 отсутствует.
\end{itemize}
\subsection{Аристотель и внешние справочные источники}

\begin{itemize}
\item Челпанов Г. И. \emph{Учебник логики (для гимназий и самообразования)}. 5-е изд. М.: Типо-литография Т-ва И. Н. Кушнерев и К°, 1909. Визуально сверены печатные с. 75–78, 97–98, 107–108 и 144 по \href{https://web.nlu.org.ua/view.html?id=238}{скану НБУ им. Ярослава Мудрого}; это выборочные страницы, не полная транскрипция издания.
\item Lewis Carroll. \emph{The Game of Logic}. London: Macmillan, 1886, печатные с. 4–5 и 9. Проверены описание двухпризнаковой диаграммы и правила для занятых, пустых и неопределённо локализованных областей по \href{https://archive.org/details/gameoflogic00carrrich}{скану Internet Archive}.
\item Н. П. Брусенцов. «Алгебраизация безупречной логики Аристотеля путем усовершенствования кэрролова метода индексов», в: Н. П. Брусенцов, Ю. С. Владимирова, \emph{Аристотелева силлогистика в символической логике Льюиса Кэрролла}. М.: Фонд «Новое тысячелетие», 2011, печатные с. 3–5. В корпусе доступен только фрагмент (PDF с. 1–3); заглавие и авторство фрагмента проверены по изображению.
\item C. I. Lewis. \emph{A Survey of Symbolic Logic}. Berkeley: University of California Press, 1918, гл. V «The System of Strict Implication», печатные с. 293 и 324. \href{https://archive.org/details/surveyofsymbolic00lewiiala}{Оцифрованный экземпляр Internet Archive}; страница с определением проверена по изображению.
\item Аристотель. \href{https://classics.mit.edu/Aristotle/prior.1.i.html}{«Первая аналитика»}, I.1.
\item Аристотель. \href{https://classics.mit.edu/Aristotle/posterior.1.i.html}{«Вторая аналитика»}, I.2.
\item Аристотель. \href{https://aristotle.johnhboyer.com/Meta/book/4/}{«Метафизика»}, Γ.7.
\item Аристотель. \href{https://www.st-andrews.ac.uk/~mnat/~ball0888/salamis/interpretatione.html}{«Об истолковании»}, гл. 9.
\item Stanford Encyclopedia of Philosophy. \href{https://plato.stanford.edu/entries/aristotle-logic/}{Aristotle’s Logic}, \href{https://plato.stanford.edu/entries/logical-consequence/}{Logical Consequence}, \href{https://plato.stanford.edu/entries/logic-relevance/}{Relevance Logic}, \href{https://plato.stanford.edu/entries/lewis-ci/}{Clarence Irving Lewis}, \href{https://plato.stanford.edu/entries/logic-conditionals/}{The Logic of Conditionals}.
\item Stanford Encyclopedia of Philosophy. \href{https://plato.stanford.edu/archives/sum2025/entries/logic-ancient/}{Ancient Logic}, раздел 2.4.
\item Ян Лукасевич. «O logice trójwartościowej», \emph{Ruch Filozoficzny}, т. 5, № 9, 1920, с. 170–171. \href{https://kpbc.umk.pl/Content/247342/PDF/EOD_OPEN_006_12_nr_09.pdf}{Скан выпуска № 9}; значения 0, 1 и 2 проверены по изображениям печатных страниц.
\item Stanford Encyclopedia of Philosophy. \href{https://plato.stanford.edu/entries/lukasiewicz/}{Jan Łukasiewicz}, § 5.1; приведена современная нормировка третьего значения как \(\tfrac12\), отличная от оригинальной записи 1920 года.
\item Patrick Blackburn and Johan van Benthem. \href{https://doi.org/10.1016/S1570-2464(07)80004-8}{“Modal Logic: A Semantic Perspective”}, in \emph{Handbook of Modal Logic}, ed. Patrick Blackburn, Johan van Benthem, and Frank Wolter, 2007, section 2.2.
\item Московский государственный университет. \href{https://ternarycomp.cs.msu.ru/Brusentsov/Brusentsov_v2.html}{Список научных трудов и изобретений Н. П. Брусенцова}; \href{https://ternarycomp.cs.msu.ru/all.html}{каталог публикаций с полными текстами}.
\end{itemize}
\section{Связанные документы проекта}

\begin{itemize}
\item \href{../docs/Саммари-идей-Брусенцова.tex}{Саммари идей Н. П. Брусенцова}.
\item \href{../docs/Как-читать-формулы-Брусенцова.tex}{Как читать формулы Н. П. Брусенцова}.
\item \href{Проверка-идей-Брусенцова-в-Lean.tex}{Заключение о проверке формализаций идей Н. П. Брусенцова в Lean}.
\item \href{Полный-указатель-трудов-Брусенцова.md}{Полный указатель работ Н. П. Брусенцова по списку МГУ}.
\item \href{Карта-работ-Брусенцова.tex}{Карта корпуса работ Н. П. Брусенцова и доступности}.
\end{itemize}
\end{document}
